\Needspace{13\baselineskip}
\subsection{A balanced metric and a finite upper mass}
\label{sec:lower-mass-metric-geometry}

Although every Hermitian form on the threefold of \cref{thm:vanishing-lower-mass} has zero lower mass, the metric $\omega$ of \eqref{eq:lower-mass-quotient-metric} is balanced and has finite upper mass. The threefold is a primitive elliptic bundle, so the loss of mass under a potential is a vertical Dirichlet energy along the elliptic fibers, and the degeneration cannot be realized by potentials that are the sum of a periodic function of $t$ and a function on the base: their masses stay above $8$. Put $D=\mathrm{d}t\wedge\dc t$, so that $\omega=\alpha+\beta+D$.

Let $B=\mathbb P^1\times\mathbb P^1$, let $\alpha_B,\beta_B$ be the Fubini--Study forms of the two factors, normalized to have integral one, so that $\alpha,\beta$ are their pullbacks, and put $\eta=\alpha_B+\beta_B$. On $B$ consider the line bundle $\mathcal O(1)\boxtimes\mathcal O(-1)$, whose fiber at $([z],[w])$ is the space of linear maps from the line $\mathbb Cz$ to the line $\mathbb Cw$, with the metric induced by the standard norms. Sending $[z:w]\in\Omega$ to the linear map $z\mapsto w$ identifies $\Omega$ with the complement of the zero section, the function $\log(\|w\|^2/\|z\|^2)$ with the logarithm of the squared norm, and $g$ with multiplication by $q=\mathrm e^{2\pi}$. The curvature representative of minus the first Chern class is $c=\beta_B-\alpha_B$, and $\eta\wedge c=0$. Hence $X$ is the primitive elliptic bundle of \cref{sec:primitive-elliptic-bundles} for these data, with the same function $t$, and $\omega$ is its metric $\pi^*\eta+D$. The constants of that subsection are
\[
 T=4\pi,\qquad A=2,\qquad c^2=-2\alpha_B\wedge\beta_B,\qquad \kappa=1,
\]
and the eigenvalues of $c$ relative to $\eta$ are $\pm1$, so that $M_0=1$.

\begin{corollary}\label{prop:lower-mass-balanced-metric}\index{Balanced metric!finite upper mass}\index{Elliptic fibration!vanishing lower mass}
On the threefold $X$ of \cref{thm:vanishing-lower-mass}, the Hermitian form $\omega=\alpha+\beta+D$ satisfies
\[
 \mathrm{d}(\omega^2)=0,\qquad \ddc\omega=2\alpha\wedge\beta\geq0,
 \qquad \int_X\omega^3=24\pi.
\]
For every smooth real function $u$,
\begin{equation}\label{eq:lower-mass-vertical-energy}
 \int_X(\omega+\ddc u)^3
 =24\pi-6\int_X\mathrm{d}u\wedge\dc u\wedge\alpha\wedge\beta.
\end{equation}
Consequently $v_+(\omega)=24\pi$. In particular, $X$ has the bounded mass property, while $v_-(h)=0$ for every Hermitian form $h$ on $X$.
\end{corollary}
\begin{proof}
\Cref{prop:primitive-elliptic-metric}, with $T=4\pi$, $A=2$ and $-\pi^*(c^2)=2\alpha\wedge\beta$, gives the balanced condition, the formula for $\ddc\omega$, the volume $3TA=24\pi$, the identity \eqref{eq:lower-mass-vertical-energy} and $v_+(\omega)=24\pi$. The metric independence in \cref{prop:bmp-metric} then gives the bounded mass property of $X$. The assertion about the lower mass is \cref{thm:vanishing-lower-mass}.
\end{proof}

For an admissible smooth $u$, equality in the upper-mass bound holds precisely when $u$ is constant along every elliptic fiber. Indeed, $\alpha\wedge\beta$ is a positive volume form in the base directions, so the nonnegative integral in \eqref{eq:lower-mass-vertical-energy} vanishes exactly when the derivative of $u$ along the fibers vanishes. The fibers are connected, and smooth local sections then show that $u$ is the pullback of a smooth function on $\mathbb P^1\times\mathbb P^1$. For the potentials with masses tending to zero in \cref{thm:vanishing-lower-mass}, the same integral instead tends to $4\pi$.

A function $f(t)$ descends to $X$ when $f$ is $4\pi$-periodic, since $t\circ g=t+4\pi$. Even after adding an arbitrary potential pulled back from the base, such functions leave a positive mass gap. The sharp gap follows from \cref{lem:periodic-variance-bound}, through the bounds on the first two derivatives of $f$ imposed by positivity.

\begin{corollary}\label{prop:elliptic-radial-mass}\index{Monge--Ampere mass@Monge--Amp\`ere mass!radial gap}
Let $\pi\colon X\to B=\mathbb P^1\times\mathbb P^1$ be the elliptic fibration above, and put $L=4\pi$. Consider the smooth potentials
\[
 u=f(t)+\pi^*v,\qquad \omega+\ddc u>0,
\]
where $f\colon\mathbb R\to\mathbb R$ is smooth and $L$-periodic and $v\in C^\infty(B,\R)$. The infimum of $\int_X(\omega+\ddc u)^3$ over these potentials is $8$. Every such mass is strictly greater than $8$, and the infimum is approached already with $v=0$.
\end{corollary}
\begin{proof}
Here $L=T$ and $\kappa^{-1/2}=M_0=1$. Since $T\geq2$, the definition of $F$ gives $F(1,4\pi)=1-1/(3\pi)$, so both bounds in \cref{prop:primitive-elliptic-separated} equal $3TA\bigl(1-F(1,4\pi)\bigr)=8$. That proposition therefore shows that every mass in the statement is strictly greater than $8$, and that potentials $f(t)$ with $v=0$ have masses tending to $8$. In terms of $p=f'$, the mass formula in its proof reads
\begin{equation}\label{eq:elliptic-radial-mass}
 \int_X(\omega+\ddc u)^3
 =6\left(L-\int_0^L p(s)^2\,\mathrm{d}s\right).
\end{equation}
\end{proof}

\Needspace{12\baselineskip}
\begin{corollary}\label{prop:lower-mass-no-pluriclosed}
The threefold $X$ admits neither a pluriclosed Hermitian metric nor a closed K\"ahler current.
\end{corollary}
\begin{proof}
The first Chern class of $\mathcal O(1)\boxtimes\mathcal O(-1)$ is represented by $\alpha_B-\beta_B$, which is nonzero in real cohomology, so this is part of \cref{prop:primitive-elliptic-metric}.
\end{proof}

Since a compact complex manifold carries a closed K\"ahler current exactly when it belongs to Fujiki's class $\mathcal C$, as recalled in \cref{sec:non-kahler-geometric-problems}, the threefold $X$ lies outside this class. The example shows that a balanced metric does not guarantee positive lower mass, even together with bounded upper mass. The Hopf threefold, in contrast, has positive lower mass and infinite upper mass, so the two properties are logically independent, and the products in \cref{cor:both-mass-properties-fail} show that both can fail at once. In particular, in a theorem asserting solvability for every nonzero finite non-pluripolar measure, bounded mass alone cannot replace the hypothesis of positive lower mass.

The exact equations of \cref{sec:zero-lower-mass-product-equations} on products $X_0\times Y$ with $v_-(h_0)=0$ apply to this threefold. We use the notation of that section with $Y=\mathbb P^2$: $\kappa$ is the Fubini--Study form normalized by $\int_Y\kappa^2=1$, $\mathrm{d}V$ is a smooth positive probability volume form on $X_0\times Y$, the functions $R$, $q_\alpha$ and $u_\alpha$, for $0<\alpha<1$, are defined there, and $b_H$ is the threshold of \cref{sec:individual-density-subsolutions}.

\Needspace{6\baselineskip}
\begin{corollary}\label{cor:elliptic-fivefold-exact-threshold}\index{Elliptic fibration!exact product equation}
Let $X_0$ be the elliptic threefold of \cref{thm:vanishing-lower-mass}, and let $h_0$ be the balanced metric of \cref{prop:lower-mass-balanced-metric}. On $X_0\times\mathbb P^2$, with background $H=h_0+\kappa$, the densities $f=Rq_\alpha$ of \cref{thm:zero-lower-mass-product-equation} satisfy
\[
 v_-(H)=0,\qquad b_H(f)=240\pi,
 \qquad\mathrm{MA}^{\mathrm{ext}}_H(u_\alpha)=240\pi f\,\mathrm{d}V.
\]
\end{corollary}
\begin{proof}
\Cref{thm:vanishing-lower-mass} gives $v_-(h_0)=0$, and \cref{prop:lower-mass-balanced-metric} gives $V_0=24\pi$. \Cref{thm:zero-lower-mass-product-equation} applies with $m=3$, $k=2$ and $C=\binom53=10$.
\end{proof}
